% Section IX: Conclusion and Future Work. Scope-B rewrite.

\subsection{Summary}
\label{sec:conclusion:summary}

We have presented a closed-loop deployment of Lexicographic Constraint Programming for trajectory Model Predictive Control on the \texttt{nuPlan} simulator. Building on the geometric equivalence framework of~\cite{lcp2025}, we have (i) released a dynamics-agnostic reference implementation that reproduces the framework's published numerical answers to $10^{-6}$ precision via a 55-test regression suite covering Algorithms~0, 1A, 1B, the pre-flight FM~I/FM~II diagnostics, the online cascade-fallback policy, the Relaxation Decision Framework, and the upper-image reification of \cite{lcp2025} \S3; (ii) realised the calibrated weighted-sum operational mode and the lex cascade as a reference inside a two-level MPC pipeline with a Sequential Linear Programming outer loop, consolidated as a roster of eleven implementation primitives plus the observer subsystem (Table~\ref{tab:lessons_learned}); and (iii) demonstrated the deployment on the 16-scenario \texttt{nuPlan-mini} benchmark in a dual-condition protocol where the operational planner ($C_1$) and the formal cascade ($C_4$) were run on each scenario.

The empirical evidence supports the four deployment claims summarised in Section~\ref{sec:results:summary}: priority-respecting compromise is realised on every scenario (the low-priority comfort and headway levels $L_0$--$L_3$ absorb the violation mass; the high-priority collision levels $L_8$--$L_{10}$ remain effectively held; Figs.~\ref{fig:top_violation}--\ref{fig:level_breakdown}); the operational weighted-sum solve reproduces the formal cascade's per-tick compliance pattern at the effectively-held levels with median match rate $1.000$ at $L_{10}$ and $L_2$, $0.987$ at $L_9$, $0.981$ at $L_8$ --- the operational realisation of \cite{lcp2025}'s Theorem~4.1 in closed loop (Fig.~\ref{fig:compliance_match}); the architectural wall-time benefit of substituting the calibrated weighted-sum solve for the $L{+}1$-stage cascade is realised at median $\rho^{(s)} = 6.22$ across the 16 scenarios (range $1.76$--$9.39$; Table~\ref{tab:walltime}); and the framework's standing assumptions (FM~I LICQ + FM~II convexity) hold across the benchmark (Table~\ref{tab:diagnostics}), with the empirical \S10.2 necessity scan (Table~\ref{tab:necessity}) identifying the scenarios on which the deployment legitimately operates in the necessity-required relaxation regime.

The present work does not address the multi-tick state-machine rules of Table~\ref{tab:rulebook} (mandatory-stop approach, yield-priority arbitration, block-the-box detection, uncontrolled-intersection negotiation, unmarked-crosswalk yield); the nine observer-only rules fall outside the convex per-step template and are not encoded by the planner. Extending the framework to encompass them is a hybrid-MPC research direction (point~3 of Section~\ref{sec:conclusion:future} below) beyond the present scope.

\subsection{Future Work}
\label{sec:conclusion:future}

For future work, we are interested in extending the deployment along four directions.

\emph{An inverse problem on the equivalence region.} The equivalence region $\Omega(p^{\star})$ of Theorem~\ref{thm:equivalence} is a polyhedron in weight space whose interior contains all weights that recover the lex optimum. The inverse problem treats $\Omega(p^{\star})$ as a feasibility set for weight inference from expert demonstrations: given a corpus of human-driven trajectories, find the weight vector $w^{\dagger}$ that best explains the demonstrations subject to lying in $\Omega(p^{\star})$. The inverse-problem formulation reconciles the framework's structural priorities with operator-supplied or learned weight preferences.

\emph{Online recalibration on first compliance-check mismatch.} The runtime compliance checker (\S\ref{sec:deployment:compliance}) currently logs mismatches but the recalibration policy is operator-supplied. The \texttt{lcp.deploy\_tick} entry point implements the \cite{lcp2025} \S9.6 cascade-fallback and Algorithm~1A/1B recalibration policies; wiring these into the production planner so the cache updates in place on the first mismatch would close the loop between the offline and online phases of the operational workflow of Fig.~\ref{fig:state_machine}. The engineering challenge is to perform the recalibration within a few ticks; the algorithmic challenge is to handle the race condition between the recalibration and the next tick's deployment.

\emph{Hybrid Model Predictive Control for state-machine rules.} The 9 observer-only rules of Table~\ref{tab:rulebook} sit outside the convex template because their semantics require discrete-mode reasoning. A hybrid extension that augments the convex per-step OCP with discrete mode variables (mandatory-stop approach states, yield-priority modes, intersection-blocking detection) would extend the lex-dominance guarantee from the 13 MPC-controlled rules to a larger subset of the rulebook at the cost of mixed-integer programming overhead.

\emph{Cross-vehicle and multi-seed deployment.} All empirical claims in this manuscript are tied to a single vehicle model (Chrysler Pacifica) and a single random seed. Cross-vehicle deployment requires re-calibration of the kinematic bicycle parameters, re-tuning of the per-rule encoder thresholds, and re-population of the calibration cache for the new platform. Multi-seed stochastic-robustness evidence would complement the single-seed qualitative pattern presented here. The framework and the deployment are vehicle-agnostic and seed-agnostic in structure; the engineering work to scale to a second platform or to a multi-seed protocol is well-scoped by the lessons roster of Table~\ref{tab:lessons_learned}.
